\subsection{Decomposing the welfare difference}
\label{ch:decomposition}

The adjoint calculation gives the full derivative but does not separate its
departure from traffic into distinct economic sources. To separate those sources,
the package compares models that share the observed baseline and policy forcing
but selectively hold congestion, modes, or routes fixed.

\subsubsection{Common-baseline comparisons}

The full primitive-cost welfare elasticity can differ from traffic for several
reasons. To isolate them, the package evaluates each alternative at the same
observed baseline and under the same policy forcing, following the
inverse-Jacobian logic in \citet{FuchsWong2026Multimodal}. Recalibrating each
alternative would mix the effect under study with a change in starting values.

The package constructs six cases:
\begin{longtable}{L{0.18\textwidth}L{0.18\textwidth}L{0.22\textwidth}L{0.27\textwidth}}
\toprule
Closure & Routes & Modes & Congestion \\
\midrule
\endhead
\code{H}: Hulten & Not solved & Not solved & Traditional approach \\
\code{NC}: no congestion & Flexible & Flexible & None \\
\code{NT}: no terminals & Flexible & Flexible & Edge congestion \\
\code{F}: full & Flexible & Flexible & All specified congestion \\
\code{FM}: fixed modes & Flexible & Fixed at baseline shares & All specified congestion \\
\code{FR}: fixed routes & Baseline OD-edge incidence fixed & Flexible & All specified congestion \\
\bottomrule
\end{longtable}

The labels describe local derivatives, not finite counterfactual models. In
particular, \code{FR} holds the differentiated route incidence fixed while
retaining modal response and the remaining equilibrium system. Both spatial
models use these comparisons. The additional normalized and
allocation--scarcity--equilibrium decompositions below are available only for
the economic-geography model.

\subsubsection{Economic-geography normalized multipliers}

To compare equilibrium propagation across cases, divide out traffic and
the common scale \(\rho\). For a positive-traffic edge \(e\) with
\(\rho\neq0\), define
\begin{equation}
 m_{e,S}=\frac{\Ereal_{e,S}}{\rho\Xi_e}.
\label{eq:guide-normalized-multiplier}
\end{equation}
The resulting multiplier measures equilibrium propagation per unit of
traffic. Define the signed wedges
\[
\begin{aligned}
 d_{edge}&=m_{NC}-m_{NT},&
 d_{terminal}&=m_{NT}-m_F,\\
 d_{mode}&=m_F-m_{FM},&
 d_{route}&=m_F-m_{FR}.
\end{aligned}
\]
Evaluating them at the common baseline yields the closure ladder
\begin{equation}
 m_F=m_{NC}-d_{edge}-d_{terminal}
    =m_{FM}+d_{mode}
    =m_{FR}+d_{route}.
\label{eq:guide-closure-ladder}
\end{equation}
The terms are signed. A positive road wedge in multiplier units need not imply
a positive contribution to the welfare elasticity when \(\rho<0\).

The first equality gives an additive congestion ladder. The other two compare
the full model separately with fixed modes and fixed routes. Mode and route
wedges are not added to the congestion ladder because each is an alternative
closure comparison.

The urban model reports the corresponding differences directly in welfare
elasticity levels:
\[
 E_{U,F}^r-E_{U,NC}^r,\qquad
 E_{U,F}^r-E_{U,FM}^r,\qquad
 E_{U,F}^r-E_{U,FR}^r.
\]
It does not divide by the economic-geography coefficient \(\rho\). The same
comparisons therefore apply without imposing a trade-model
normalization on residence--workplace welfare.

\subsubsection{Inverse-gap identity}

Each rung compares two derivatives that share the welfare row and policy
forcing but use different Jacobians. The matrix identity
\(J_A^{-1}-J_B^{-1}=J_A^{-1}(J_B-J_A)J_B^{-1}\) therefore implies
\[
 q_c^\top J_{c,A}^{-1}b_c-q_c^\top J_{c,B}^{-1}b_c
 =q_c^\top J_{c,A}^{-1}
 (J_{c,B}-J_{c,A})J_{c,B}^{-1}b_c.
\]
The package checks this identity for the edge, terminal, mode, and route
comparisons in both spatial models. It is a direct algebraic test of the
signs, orientation, and Jacobian construction.

\subsubsection{Economic-geography analytical components}

The inverse-gap identity gives the total difference between two closures. For
economic geography, the package further separates each mode or route wedge
into allocation, congestion-scarcity, and aggregate-equilibrium components:
\[
 d_r=d_r^{alloc}+d_r^{scar}+d_r^{eq},
 \qquad r\in\{mode,route\}.
\]
Under the common-baseline construction used here, the sparse spatial
allocation block and aggregate-feedback term are the same across the five
closures.
Their contributions are therefore exact zeros for the road, terminal, mode,
and route comparisons, and the congestion-scarcity component carries each
wedge. The package nevertheless evaluates the general mixed update for the
mode and route comparisons and checks both the implied zeros and the
reconstruction of the total wedge. Appendix \ref{app:algorithm} gives the
Jacobian blocks, rank-two update, and structured inverse formulas.

\subsubsection{The economic-geography Hulten gap}

Substituting
\(\Ereal_{e,F}=\rho\Xi_e m_{e,F}\) from
\eqref{eq:guide-normalized-multiplier} into the difference between traffic and
the full realized derivative gives
\[
 \Xi_e-\Ereal_{e,F}
 =\Xi_e(1-\rho)+\rho\Xi_e(1-m_{e,F}).
\]
The first term is the common externality scaling. The second is equilibrium
propagation relative to the Traditional approach.

For a nonzero realized derivative, define
\(\chi_e^{eff}=\Eprimitive_{e,F}/\Ereal_{e,F}\). Next substitute the first equality in
\eqref{eq:guide-closure-ladder} and use
\(\Eprimitive_{e,F}=\chi_e^{eff}\Ereal_{e,F}\). This produces
\begin{align*}
 \Xi_e-\Eprimitive_{e,F}
 ={}&\Xi_e(1-\rho)
 +\rho\Xi_e(1-m_{e,NC})\\
 &+\rho\Xi_e d_{edge,e}
 +\rho\Xi_e d_{terminal,e}
 +\rho\Xi_e m_{e,F}(1-\chi_e^{eff}).
\end{align*}
A zero realized derivative leads the output to compute the
primitive-minus-realized contribution directly in levels instead of forming a
ratio. Signed percentages can be misleading if large positive and negative
components nearly cancel.

The identity reports all terms in units of the Traditional approach minus
primitive-cost welfare. A positive contribution widens that gap; a negative
contribution offsets it. The sign convention applies even when \(\rho<0\).

For the urban model, the package reports
\(\Xi_e-E_{U,F}^r\) and
\(\Xi_e-E_{U,F}^{\theta}\), together with the closure gaps above. It does not
label these differences as externality scaling or trade-market-access
propagation because those labels come from the economic-geography
factorization.

\subsubsection{Physical-link aggregation}

The model is directed. To obtain a two-direction physical-link experiment, the
software first sums the primitive and realized elasticities across the two
directions. It applies the same operation to Hulten traffic, all additive
Hulten-gap components, and all analytical components.
Only after summation does it calculate a physical-link normalized multiplier.
This order avoids averaging ratios with different traffic denominators.

\warningbox{A physical-link result is defined only when exactly two declared
policy arcs have opposite endpoints and a common physical-link identifier.
Transit services and one-way links that lack such a pair remain directed
policies.}

\subsubsection{Reading a decomposition}

Begin with the gap between the Traditional approach and the full
primitive-cost result. In the economic-geography model, separate the common
externality scale from link-specific equilibrium propagation before examining
edge and terminal congestion. Next compare the primitive- and realized-cost
effects to measure pass-through. The fixed-mode and fixed-route cases then show
whether modal or route adjustment contributes to ranking changes. Interpret
these economic differences only after the identity residuals have passed their
numerical checks.
The decomposition explains a result within the specified model. It does not
show which parameters or model assumptions are empirically appropriate.

Applying these identities requires a baseline whose flows, units, and policy
definition satisfy the model's accounting restrictions. Subsection
\ref{ch:data} states that empirical contract before the guide turns to the
software.
